\documentclass[11pt]{article}
\usepackage[T1]{fontenc}
\usepackage[a4paper,margin=1in]{geometry}
\usepackage{amsmath,amssymb,amsthm}
\usepackage[hidelinks]{hyperref}
\usepackage{microtype}

\newtheorem*{question}{Source question}
\newtheorem*{answer}{Literature answer}

\setlength{\parindent}{0pt}
\setlength{\parskip}{0.6em}

\title{Literature Status: Endpoint Lorentz Estimates\\
for All Cocancelling Constraints}
\author{Run \texttt{fa\_banach\_001} --- lane 2}
\date{11 August 2026}

\begin{document}
\maketitle

\begin{center}
\fbox{\parbox{0.9\textwidth}{
\textbf{Status: literature already answered.}
The broad cocancelling-operator question in arXiv:2008.05639 is answered
affirmatively.  Stolyarov's arXiv:2010.05297 proves a stronger
Besov--Lorentz theorem for every closed translation- and
dilation-invariant constrained space excluding point masses.  For
$\mathcal W=\ker L(D)$, exclusion of nonzero $a\otimes\delta_0$ is exactly
cocancellation.  Breit--Cianchi--Spector, arXiv:2512.06352, state the
result directly for arbitrary homogeneous cocancelling operators and
attribute the first- and higher-order cases to the intervening literature.
}}
\end{center}

\section{Original question}

F. Hernandez and D. Spector,
\emph{Fractional Integration and Optimal Estimates for Elliptic Systems},
arXiv:2008.05639, later published in \emph{Calculus of Variations and
Partial Differential Equations} \textbf{63} (2024), Paper No.~117
\cite{source}, prove the endpoint estimate for divergence-free fields.
In Section~2, arXiv PDF page~9, they isolate the broader question:

\begin{question}
Does the optimal endpoint Lorentz estimate hold for every cocancelling
operator?  Equivalently, if $L(D)$ is a homogeneous cocancelling
differential operator, $0<\alpha<d$, and $L(D)F=0$, does one have
\[
  \|I_\alpha F\|_{L^{d/(d-\alpha),1}(\mathbb R^d)}
  \leq C\|F\|_{L^1(\mathbb R^d)}?
\]
\end{question}

The source says explicitly that this was not known and that an affirmative
answer would resolve the open problems listed immediately beforehand.

\section{Answer and exact identification}

D. Stolyarov, \emph{Hardy--Littlewood--Sobolev inequality for $p=1$},
arXiv:2010.05297v3, later published in \emph{Sbornik: Mathematics}
\textbf{213} (2022), 844--889 \cite{stolyarov}, proves the stronger
Besov--Lorentz statement.  His Theorem~1.2 (arXiv PDF page~7) says that if
$\mathcal W\subset\mathcal S'(\mathbb R^d;\mathbb R^\ell)$ is closed and
translation- and dilation-invariant, then, for $0<\alpha<d$,
\[
 I_\alpha:\mathcal W\cap L^1
 \longrightarrow \dot B^{0,1}_{d/(d-\alpha),1}
\]
is bounded if and only if $\mathcal W$ contains no nonzero
$a\otimes\delta_0$.  Equation~(1.18) and Corollary~1.7 (arXiv PDF page~3)
give the Lorentz consequence
\[
  \|I_\alpha F\|_{L^{d/(d-\alpha),1}}
  \lesssim \|F\|_{L^1}.
\]

Take $\mathcal W=\{F:L(D)F=0\}$.  This space is closed and invariant under
translations and dilations because $L(D)$ is homogeneous with constant
coefficients.  Moreover,
\[
 L(D)(a\delta_0)=0
 \quad\Longleftrightarrow\quad
 L(\xi)a=0\ \text{for every }\xi\ne0
 \quad\Longleftrightarrow\quad
 a\in\bigcap_{\xi\ne0}\ker L(\xi).
\]
Thus $\mathcal W$ contains no nonzero point mass exactly when $L$ is
cocancelling.  Stolyarov's theorem therefore answers the source question
for operators of arbitrary homogeneous order, with a stronger target norm.

This identification is also explicit in later literature.  On arXiv PDF
page~3 of Breit--Cianchi--Spector,
\emph{Riesz potential estimates under co-canceling constraints},
arXiv:2512.06352 \cite{bcs}, equation~(1.8) states the desired
$L^1\to L^{n/(n-\alpha),1}$ estimate for every homogeneous cocancelling
operator of order $k\in\mathbb N$.  The following sentence attributes the
first-order case to Hernandez--Rai\c{t}\u{a}--Spector and the higher-order
case to Stolyarov.  Their Theorem~3.1 (arXiv PDF page~10) further extends
the statement to general rearrangement-invariant domain and target norms.

\section{Scope and search status}

The all-orders Euclidean cocancelling question on arXiv PDF page~9 is
completely answered affirmatively.  This status does \emph{not} by itself
settle the distinct periodic-torus endpoint question mentioned on source
PDF page~4.

A bounded search through 11 August 2026 checked the run registry, the local
arXiv corpus, the exact question wording, the phrases ``cocancelling
operator'' and ``endpoint Lorentz'', and later papers citing the source.
The decisive sources are arXiv:2010.05297 and arXiv:2512.06352.  No new
mathematical result is claimed in this packet.

\begin{thebibliography}{9}
\bibitem{source}
F. Hernandez and D. Spector,
\emph{Fractional Integration and Optimal Estimates for Elliptic Systems},
Calc. Var. Partial Differential Equations \textbf{63} (2024), Paper No.~117;
arXiv:2008.05639.

\bibitem{stolyarov}
D. Stolyarov,
\emph{Hardy--Littlewood--Sobolev inequality for $p=1$},
Sb. Math. \textbf{213} (2022), 844--889;
arXiv:2010.05297.

\bibitem{bcs}
D. Breit, A. Cianchi, and D. Spector,
\emph{Riesz potential estimates under co-canceling constraints},
arXiv:2512.06352 (2025).
\end{thebibliography}

\end{document}
